\documentclass[11pt,a4paper]{article}

\usepackage[T1]{fontenc}
\usepackage[utf8]{inputenc}
\usepackage[french]{babel}
\usepackage{geometry}
\usepackage{enumitem}
\usepackage{hyperref}
\usepackage{titlesec}
\usepackage{tabularx}
\usepackage{array}
\usepackage{colortbl}
\usepackage{xcolor}

\geometry{margin=2.2cm}
\setlength{\parskip}{0.35em}
\setlist[itemize]{noitemsep, topsep=2pt, leftmargin=1.2em}
\setlist[enumerate]{noitemsep, topsep=2pt, leftmargin=1.5em}

\definecolor{ifriblue}{RGB}{45,55,120}
\definecolor{ifrigray}{RGB}{245,246,250}

\titleformat{\section}{\color{ifriblue}\normalfont\bfseries\large}{\thesection.}{0.5em}{}
\titleformat{\subsection}{\normalfont\bfseries\normalsize}{\thesubsection}{0.5em}{}

\hypersetup{colorlinks=true, linkcolor=ifriblue, urlcolor=ifriblue}

\newcommand{\metarow}[2]{#1 & #2 \\}

\begin{document}

% ===== PAGE DE GARDE =====
\begin{center}
    \colorbox{ifriblue}{\parbox{\dimexpr\textwidth-2\fboxsep}{
        \centering
        \color{white}
        \vspace{0.5em}
        {\LARGE \textbf{Rapport d'exercice --- V2}}\\[0.35em]
        {\large Gestion de Projet Informatique}
        \vspace{0.5em}
    }}
\end{center}

\vspace{0.8em}

\noindent
\begin{tabularx}{\textwidth}{|>{\bfseries}l|X|}
    \hline
    \rowcolor{ifrigray}
    \metarow{Étudiant}{Juste TOSSOU}
    \metarow{Filière}{M1-GL / SIRI / SI}
    \metarow{Établissement}{IFRI --- Institut de Formation et de Recherche en Informatique}
    \metarow{Enseignant}{Dr. Arnaud AHOUANDJINOU}
    \metarow{Année académique}{2024--2025}
    \metarow{Objectif}{Rappels du cours}
    \metarow{Consigne}{Rédiger un résumé de 2 à 3 pages maximum des notions essentielles vues en cours}
    \metarow{Données}{Support de cours (PMBOK\textsuperscript{\textregistered} 6\textsuperscript{e} et 7\textsuperscript{e} éditions, PMI)}
    \hline
\end{tabularx}

\vspace{0.8em}
\hrule
\vspace{0.6em}

% ===== CORPS =====

\section{Fondamentaux : projet, programme et portefeuille}

Un \textbf{projet} est une initiative \textbf{temporaire} orientée vers un \textbf{résultat unique}. Il vise un objectif précis avec des ressources limitées et un échéancier défini. On le distingue d'une activité opérationnelle par son caractère non répétitif, son unicité et sa capacité à induire le changement.

Le \textbf{management de projet} regroupe l'ensemble des méthodes permettant de \textbf{planifier}, \textbf{estimer} et \textbf{piloter} les activités pour livrer dans les délais, le budget et le périmètre convenus. Le travail est structuré en \textbf{phases} produisant des \textbf{livrables}, l'ensemble formant le \textbf{cycle de vie}.

\begin{itemize}
    \item \textbf{Programme} : coordination de projets liés pour maximiser les synergies.
    \item \textbf{Portefeuille} : alignement des projets/programmes sur la stratégie de l'organisation.
    \item \textbf{Porte de phase (GO/NO-GO)} : point de décision en fin de phase (continuer, ajuster ou arrêter).
\end{itemize}

\section{Acteurs clés et documents de référence}

\textbf{Acteurs principaux} : sponsor, chef de projet, équipe projet, comité de pilotage (COPIL), PMO, clients, utilisateurs finaux, fournisseurs et parties prenantes.

\textbf{Documents incontournables} :
\begin{itemize}
    \item \textbf{Business case} : viabilité économique (TIR, VAN, ROI).
    \item \textbf{Charte de projet} : autorisation officielle, objectifs, sponsor, autorité du chef de projet.
    \item \textbf{Énoncé de périmètre} : description détaillée et base d'accord entre parties prenantes.
    \item \textbf{Exigences} : besoins formalisés, quantifiés et traçables.
\end{itemize}

\section{Le chef de projet : rôle, leadership et équipe}

Le chef de projet est un \textbf{coordinateur} : il assure la vision globale, la communication (environ 90\,\% de son temps) et l'atteinte des objectifs sans exécuter seul toutes les tâches métiers. Il arbitre le triangle \textbf{coût -- délai -- périmètre/qualité}.

\textbf{Leadership vs management} : inspirer et mobiliser (leadership) ; organiser et contrôler (management).

Outils et référentiels utiles :
\begin{itemize}
    \item Objectifs \textbf{SMART} (Spécifique, Mesurable, Atteignable, Réaliste, Temporellement défini).
    \item \textbf{Charte d'équipe} et \textbf{règles de base} (ground rules).
    \item Motivation : pyramide de \textbf{Maslow}, deux facteurs de \textbf{Herzberg}.
    \item \textbf{Servant leader} (Agile) : répondre aux besoins de l'équipe et du client.
\end{itemize}

\section{Gestion des conflits}

Les conflits sont fréquents en projet (ressources, priorités, méthodes, cultures). Le chef de projet doit les anticiper et les traiter constructivement.

\textbf{Origines} : concurrence, désaccords sur les rôles, divergences d'objectifs, mauvaise communication.

\textbf{Stratégies} : évitement, accommodement, compromis, forçage, \textbf{collaboration} (solution gagnant-gagnant, recommandée).

\section{PMI et évolution PMBOK 6 $\rightarrow$ 7}

Le \textbf{PMI} (Project Management Institute) publie le \textbf{PMBOK\textsuperscript{\textregistered}} et certifie les praticiens (ex. \textbf{PMP}).

\vspace{0.3em}
\noindent
\small
\begin{tabularx}{\textwidth}{|>{\bfseries}p{2.8cm}|X|X|}
    \hline
    \rowcolor{ifrigray}
    & \textbf{PMBOK 6} & \textbf{PMBOK 7} \\
    \hline
    Logique & Processus (EOTS) & Principes et valeur \\
    Structure & 10 domaines de connaissance & 8 domaines de performance \\
    Pilotage & 5 groupes de processus & 12 principes de gestion \\
    Focus & Livrables & Résultats et valeur métier \\
    Approches & Prédictif dominant & Prédictif, Agile, Hybride \\
    \hline
\end{tabularx}
\normalsize

\vspace{0.4em}
\textbf{5 groupes de processus (V6)} : initialisation, planification, exécution, maîtrise, clôture.

\section{Prédictif, Agile, Scrum et cycles de vie}

\textbf{Prédictif (Waterfall)} : périmètre stable, planification complète en amont. Structuré mais peu flexible.

\textbf{Agile} : itérations courtes, adaptation continue. Le \textbf{Manifeste Agile} privilégie les personnes, le produit fonctionnel, la collaboration et l'accueil du changement.

\textbf{Scrum} en bref :
\begin{itemize}
    \item \textbf{Rôles} : Product Owner, Scrum Master, Développeurs.
    \item \textbf{Artéfacts} : Product Backlog, Sprint Backlog, Incrément.
    \item \textbf{Cérémonies} : Sprint Planning, Daily Scrum, Sprint Review, Sprint Retrospective.
    \item \textbf{User story} : \textit{En tant que [utilisateur], je veux [action], afin de [bénéfice]}.
\end{itemize}

\textbf{Hybride} : combinaison contextualisée (ex. planification classique + rétrospectives agiles).

\textbf{Cycles de vie possibles} : prédictif, itératif, incrémental, agile, hybride.

\section{Organisation, PMO, EEF et OPA}

\textbf{Structures} :
\begin{itemize}
    \item \textbf{Fonctionnelle} --- faible autorité du chef de projet.
    \item \textbf{Orientée projet} --- forte autorité, équipe dédiée.
    \item \textbf{Matricielle} --- compromis entre les deux (faible, équilibrée, forte).
\end{itemize}

\textbf{PMO} : bureau de projet (soutenant, conformité, directif) standardisant méthodes, outils et gouvernance.

\textbf{EEF} : facteurs environnementaux hors contrôle direct (culture, normes, infrastructure).

\textbf{OPA} : actifs organisationnels réutilisables (modèles, procédures, leçons apprises, archives).

\section{Code éthique et conduite professionnelle (PMI)}

Quatre valeurs fondamentales :
\begin{itemize}
    \item \textbf{Responsabilité} : assumer ses choix et leurs conséquences.
    \item \textbf{Respect} : considération envers soi, les autres et les ressources.
    \item \textbf{Équité} : impartialité et absence de favoritisme.
    \item \textbf{Honnêteté} : transparence et intégrité dans toutes les interactions.
\end{itemize}

\vspace{0.8em}
\hrule
\vspace{0.5em}

\noindent
\begin{flushright}
    \textbf{Juste TOSSOU}\\
    \textit{Résumé de cours --- Gestion de Projet Informatique (V2)}
\end{flushright}

\end{document}
